\section{Solutions for Chapter~\ref{ch:regularity}}

\paragraph{Exercise~\ref{ex:c-regular-1}: the graded quotient.}
\phantomsection
\label{sol:ex:c-regular-1}
For homogeneous $a,b$, the vector $a_{(-2)}b$ has weight
$\wt(a)+\wt(b)+1$. Hence $C_2(V)$ is a direct sum of homogeneous
subspaces. In finite-dimensional $V_n$ every subspace has an
orthogonal complement, and
$V_n/(C_2(V)\cap V_n)$ is represented by that complement.
This is finite-dimensional linear algebra grade by grade; no
assertion about Hilbert-space closure of the infinite direct sum
is needed.



\paragraph{Exercise~\ref{ex:c-regular-2}: the trace step.}
\phantomsection
\label{sol:ex:c-regular-2}
A diagonalizable operator with nonnegative eigenvalues has trace the
sum of those eigenvalues. Zero trace forces each to vanish, so the
operator vanishes. Without diagonalizability,
$\left(\begin{smallmatrix}0&1\\0&0\end{smallmatrix}\right)$ has
zero spectrum and trace but is nonzero. This is precisely why the
regular component top decomposition is used on a generalized
weight-zero space: nilpotence of the total energy alone is insufficient.
The argument does not require a positive form on the module carrying
these operators. It uses diagonalizability and the signs of their
eigenvalues, both supplied separately by the component theory.



\paragraph{Exercise~\ref{ex:c-regular-3}: downward induction to the vacuum condition.}
\phantomsection
\label{sol:ex:c-regular-3}
For each fixed $a$, lower truncation supplies an integer $N\geq3$ with
$a_{(n)}v=0$ for all $n\geq N$. Since $L_{-1}v=0$,
translation covariance implies
$(n+1)a_{(n)}v=-L_{-1}a_{(n+1)}v$.
Descending from $N$ first reaches $a_{(3)}v=0$, and then gives
\[
 3a_{(2)}v=-L_{-1}a_{(3)}v=0,\qquad
 2a_{(1)}v=-L_{-1}a_{(2)}v=0,\qquad
 a_{(0)}v=-L_{-1}a_{(1)}v=0.
\]
The coefficient is nonzero for every $n\geq0$. At $n=-1$ it
would vanish, exactly where a creation mode may act nontrivially.
In particular the argument imposes no erroneous vanishing of
$a_{(-1)}v$, which is the desired copy of the state $a$.



\paragraph{Exercise~\ref{ex:c-regular-4}: a nonsplit extension.}
\phantomsection
\label{sol:ex:c-regular-4}
Let $A=\C[\varepsilon]/(\varepsilon^2)$. Its regular module fits
into
$0\to\C\varepsilon\to A\to A/(\varepsilon)\to0$,
with both endpoints the simple module on which $\varepsilon$ acts
as zero. A section would send $1$ to $1+b\varepsilon$, but
$\varepsilon(1+b\varepsilon)=\varepsilon\ne0$, so no section is
$A$-linear. The simple module is therefore not projective. In the
VOA proof, the positive stress identity supplies the missing
vacuum-like lift and thereby rules out the analogous nonsplit
extension onto the unit.



\paragraph{Exercise~\ref{ex:c-regular-5}: factorial versus exponential.}
\phantomsection
\label{sol:ex:c-regular-5}
For $n$ even, the final $n/2$ factors of $n!$ are at least $n/2$,
so $n!\geq(n/2)^{n/2}$. Once $n/2>A^2$, this is greater than
$A^n$. A fixed prefactor can be absorbed by choosing $n$ larger.
To see this uniformly for even and odd $n$, put $r_n=n!/A^n$.
For $n+1\geq2A$ one has $r_{n+1}/r_n=(n+1)/A\geq2$;
hence $r_n$ eventually exceeds every fixed constant.
In the moderate-growth argument take
$A=C_X^2$ and the fixed constant $d_0(\one)^2$. Then, for all
sufficiently large $n$,
\[
 \dim\End(X^{\otimes n})
 \leq d_0(\one)^2C_X^{2n}<n!,
\]
which is exactly the hypothesis needed there.



\paragraph{Exercise~\ref{ex:c-regular-6}: the logical uses of the hypotheses.}
\phantomsection
\label{sol:ex:c-regular-6}
Positive coefficients let separate upper bounds on partial-energy
expectations add to an upper bound on ambient energy. They also
make a zero sum of nonnegative partial traces force each trace
to vanish. Component regularity supplies ordinary semisimple
restrictions, finitely many simple types, and, together with its
$C_2$ consequences, the finite component quotient bounds.
Grading compatibility makes component modes have their expected
ambient degrees. It places a lowest ambient space in complete
component tops and makes each partial energy preserve finite
ambient grades.

The unique simple component module at weight zero identifies those
zero-energy tops with vacuum tops. Their translation operators
therefore annihilate the selected vectors, giving the vacuum-like
lift needed for projectivity. If another simple had weight zero,
the trace argument would still force its energy to vanish, but
would not force its translation mode to vanish. Both uniqueness
of the $V$ module of weight zero and vacuum projectivity would
lose their stated proof. The $C_2$ bound itself would remain valid
if the other component hypotheses remained true.

Braiding is a premise of the categorical rigidity theorem used after
projectivity. Before rigidity, right exactness turns a surjection
$P^{\oplus d_0(Z)}\twoheadrightarrow Z$ into one onto $Z\otimes X$;
this gives the exponential bound on tensor-power endomorphisms.
Only after rigidity has been established is tensoring known to be
exact, allowing projectivity of the unit to imply projectivity of
every object. Interchanging these stages would be circular.


